\paragraph{\trCompactness。}\label{sub:compactness}
带\trLabel{}的\trInhomogeneous\trPseudoholomorphic\trPolygon{}之诸\trModuliSpace{}具有自然的\trGromov\trCompactification，其结构在一定程度上仿照 $\scr{R}^{d+1}\subset\bar{\scr{R}}^{d+1}$ 的结构（事实上可以把后者看成 $M=L_0=\cdots=L_d$ 为一个点时的平凡情形）。我们将把这些\trCompactification{}当作集合来定义，不加证明地叙述其主要\trTopology{}性质，然后再解释这一切如何蕴含诸映射 $\mu^d$ 构成一个\trAinfty{}结构。
按\eqref{eq:polygon-moduli-space} 的定义取定诸 $L_k$ 与 $y_k$，并取定一棵半稳定的 $d$ \trLeafedTree{}$T$。设对每个\trVertex{}$v\in\mathrm{Ve}(T)$ 都有一个 $|v|$ \trPointedDisc{}$S_v$。若 $|v|\ge 3$，该\trDisk{}决定一个点 $r_v\in\scr{R}^{|v|}$，并且存在唯一的\trIsomorphism{}$S_v\cong\scr{S}^{|v|}_{r_v}$。在剩下的情形 $|v|=2$，有 $S_v\cong Z$，但该\trIsomorphism{}只在相差 $s$ 方向\trTranslation{}的意义下确定。此外，诸\trDisk{}还须带有\trLagrange\trLabel，使得对每个\trFlag{}$f=(v,e)$ 都有一对取定的\trLagrange\trSubmanifold{}$(L_{f,0},L_{f,1})$。与前面关于\trGluing{}的讨论一样，诸\trLabel{}应当彼此\trCompatible，使得\trGluing{}而成的 $(d{+}1)$ \trPointedDisc{}上所得的\trLabel{}正是给定的诸 $L_k$。最后，对每个 $f$ 我们要有一个点 $\tilde{y}_f\in\scr{C}(L_{f,0},L_{f,1})$，满足条件 $\tilde{y}_{f^+(e)}=\tilde{y}_{f^-(e)}$，并且与\trRoot{}及诸\trLeaf{}相邻的那些\trFlag{}所关联的点按自然次序正是 $y_0,\ldots,y_d$。一个\textbf{\trInhomogeneous\trPseudoholomorphic{}稳定\trPolygon}（inhomogeneous pseudo-holomorphic stable polygon）是一族映射 $u_v\colon S_v\to M$（$v\in\mathrm{Ve}(T)$），满足下列条件：
\begin{itemize}
\item 若 $|v|\ge 3$，可以把 $u_v$ 看成映射 $\scr{S}^{|v|}_{r_v}\to M$，要求它应是相应的\trInhomogeneous\trPseudoholomorphicMap{}方程的一个解：
\begin{equation*}
(r_v,u_v)\in\scr{M}^{|v|}(\tilde{y}_{f_0(v)},\ldots,\tilde{y}_{f_{|v|-1}(v)}).
\end{equation*}
其中 $f_0(v),\ldots,f_{|v|-1}(v)$ 是按自然次序绕着 $v$ 的诸\trFlag。这是相对于从我们的\trUniversal{}选取中取来的\trPerturbationDatum{}而言的。
\item 若 $|v|=2$，则 $u_v\colon Z\to M$ 应满足
\begin{equation*}
u_v\in\scr{M}^2(\tilde{y}_{f_0(v)},\tilde{y}_{f_1(v)}),
\end{equation*}
意即：在作为我们设定之一部分的\trFloerDatum{}下，它是\trFloerEquation{}的一个非\trStationary{}解。由于该方程在\trTranslation{}下不变，这个陈述不依赖于\trIsomorphism{}$S_v\cong Z$ 的选取。
\end{itemize}
令 $\scr{M}^T(y_0,\ldots,y_d)$ 为这类\trInhomogeneous\trPseudoholomorphic{}稳定\trPolygon{}的一切\trEquivalence{}类所成的集合。该\trEquivalence{}关系是诸\trPointedDisc{}$S_v$ 的\trIsomorphism，连同诸映射 $u_v$ 之间相应的相等。具体地说，这意味着：若有一个 $|v|=2$ 的 $v$，则将相应的 $u_v$ 作一个\trTranslation{}并不影响点 $\{u_v\}\in\scr{M}^T(y_0,\ldots,y_d)$。作为集合，\trGromov\trCompactification{}定义为对一切 $T$ 取并：
\begin{equation}\label{eq:polygon-compactification}
\bar{\scr{M}}^{d+1}(y_0,\ldots,y_d)=\coprod_T\scr{M}^T(y_0,\ldots,y_d).
\end{equation}
与\eqref{eq:dm-compactification} 不同，半稳定的 $T$ 有无穷多个，因为总可以再插入更多\trValency{}为 $2$ 的\trVertex。然而，对 $(y_0,\ldots,y_d)$ 的任何取法，\eqref{eq:polygon-compactification} 中实际上非空的\trStratum{}只有有限多个。为看清这一点，把各分支的\trEnergy{} \eqref{eq:energy-definition} 加起来便得
\begin{equation*}
\sum_v E(u_v)=A_{H_{L_0,L_d}}(y_0)-\sum_{k=1}^{d}A_{H_{L_{k-1},L_k}}(y_k)+(\text{\trCurvature{}项}).
\end{equation*}
\trCurvature{}项的贡献只来自 $|v|\ge 3$ 的诸\trVertex；对我们的\trPerturbationDatum\trUniversal{}选取考察其\trConsistency{}条件之后，即可断言它在整个空间\eqref{eq:polygon-compactification} 上有一个\trUniform{}的界。另一方面，存在 $\epsilon>0$，使得每个 $|v|=2$ 的分支 $u_v$ 对\eqref{eq:energy-definition} 左端的贡献至少是 $\epsilon$。具体而言：考察所有的元素对 $y_0,y_1\in\scr{C}(L_i,L_j)$（$0\le i<j\le d$），并把 $\epsilon$ 定义为所有正的作用差 $A_{H_{L_i,L_j}}(y_0)-A_{H_{L_i,L_j}}(y_1)$ 中的最小者。这两个事实合起来，就对那些对\eqref{eq:polygon-compactification} 贡献非空\trStratum{}的\trTree，给出了其\trValency{}为 $2$ 的\trVertex{}个数的一个界。
从现在起，假设所涉及的一切\trModuliSpace{}都是\trRegular{}的。以 $\scr{M}^{d+1}(y_0,\ldots,y_d)^p$ 表示 $\scr{M}^{d+1}(y_0,\ldots,y_d)$ 的 $p$ 维部分（它是该空间的一个\trOpen\trClosed{}子集），并类似地记 $\scr{M}^T(y_0,\ldots,y_d)^p$。\eqref{eq:polygon-compactification} 上有一个自然的\trStableMap\trTopology，相对于它，整个空间是\trCompact{}的。这个\trTopology{}使得：每一个
\begin{equation}\label{eq:stratum-topology}
\bar{\scr{M}}^{d+1}(y_0,\ldots,y_d)^p=\coprod_T\scr{M}^T(y_0,\ldots,y_d)^{p-|\mathrm{Ed}^{\mathrm{int}}(T)|}
\end{equation}
都是\trOpen\trClosed{}子集，并且是\trCorneredBoundary{}的 $p$ 维\trTopology\trManifold。更确切地说，与单\trVertex{}\trTree{}相对应的部分——它可与 $\scr{M}^{d+1}(y_0,\ldots,y_d)^p$ 等同——是内部，而其余部分都是\trCodimension{}为 $|\mathrm{Ed}^{\mathrm{int}}(T)|$ 的\trBoundaryStratum。关于\eqref{eq:stratum-topology} 之可微结构的问题相当微妙：每个\trStratum{}无疑是一个\trSmooth\trManifold，而且可以定义一些\trGluingMap，给出诸\trBoundaryStratum{}的\trCollar{}邻域。然而，\trGluingParameter{}的选取就不那么清楚了，这意味着这些映射未必能像 $\bar{\scr{R}}^{d+1}$ 的情形那样互相吻合。所幸，\trFukayaCategory{}的构造只用到维数 $p\le 1$ 的诸\trModuliSpace，对它们而言这个问题无关紧要。
先考虑 $p=0$ 的情形，可见只有单\trVertex{}\trTree{}对\eqref{eq:stratum-topology} 有贡献。该空间的\trCompactness{}意味着 $\scr{M}^{d+1}(y_0,\ldots,y_d)^0$ 是有限集，这就为定义\eqref{eq:higher-compositions} 中的计点做法提供了依据。其次，对 $p=1$，边界点来自恰有两个\trVertex{}的半稳定\trTree，即\cref{fig:two-vertex-tree} 中的 $T=T(m,n)$。令
\begin{equation}\label{eq:two-vertex-boundary}
\begin{aligned}
(r_1,u_1,r_2,u_2)\in{}&\scr{M}^{d+2-m}(y_0,y_1,\ldots,y_n,\tilde{y},y_{n+m+1},\ldots,y_d)^0\\
&\times\scr{M}^{m+1}(\tilde{y},y_{n+1},\ldots,y_{n+m})^0
\end{aligned}
\end{equation}
是 $T(m,n)$ \trStratum{}中的一个点，这意味着 $\tilde{y}\in\scr{C}(L_n,L_{n+m})$。这类 $(r_1,u_1,r_2,u_2)$ 模 2 的个数——对一切可能的 $\tilde{y}$ 取并——恰好就是下式中 $y_0$ 的\trCoefficients：
\begin{equation*}
\mu^{d+1-m}(y_d,\ldots,y_{n+m+1},\mu^m(y_{n+m},\ldots,y_{n+1}),y_n,\ldots,y_1).
\end{equation*}
对 $m$ 与 $n$ 求和便把所有边界点都考虑了进来，故必得一个偶数；这正是带 $\KK$-\trCoefficients{}的\trAinfty{}结合性方程\eqref{eq:aoo-associativity} 所说的内容。
\begin{figure}[htbp]
\centering
\begin{tikzpicture}[line cap=round,line join=round,scale=1.0]
  \draw (-2.6,0)--(0,0);
  \draw (0,0)--(1.05,2.05);
  \draw (0,0)--(3.35,0.95);
  \draw (0,0)--(2.1,-0.75);
  \draw (0,0)--(1.5,-1.55);
  \draw (3.35,0.95)--(3.95,2.0);
  \draw (3.35,0.95)--(4.85,1.1);
  \draw (3.35,0.95)--(4.5,0.35);
  \node at (-0.3,0.32) {$v_1$};
  \node at (3.3,1.32) {$v_2$};
  \node at (2.5,0.3) {第 $(n{+}1)$ 条\trExteriorEdge};
  \node at (-1.15,-0.75) {$|v_1|=d+2-m$};
  \node at (-1.0,-1.25) {$|v_2|=m+1$};
\end{tikzpicture}
\caption{}\label{fig:two-vertex-tree}
\end{figure}
从一个更贴近直观的视角来考察\trGromov\trCompactification{}中最后提到的这种情形，是富有教益的。设有一列 $(r^i,u^i)\in\scr{M}^{d+1}(y_0,\ldots,y_d)^1$，\trConvergence{}到某个边界点\eqref{eq:two-vertex-boundary}。这涵盖两种相互关联而又彼此不同的现象：
\begin{itemize}
\item \emph{非稳定情形}出现于 $m=1$ 或 $m=d$ 时。为使记号简单，我们只讨论前一种情形。诸 $u^i$ 是映射 $\scr{S}^{d+1}_{r^i}\to M$，而 $u_1$ 是映射 $\scr{S}^{d+1}_{r_1}\to M$，其中 $r^i,r_1\in\scr{R}^{d+1}$。事实上，$r^i$ \trConvergence{}到 $r_1$，而诸 $u^i$ 在紧子集上\trConvergence{}到 $u_1$。另一个分支 $u_2\in\scr{M}^1(\tilde{y},y_{n+1})$\footnote{译注：原书此处印作 $\scr{M}^1$；疑应为 $\scr{M}^2$（按\eqref{eq:two-vertex-boundary} 第二因子 $\scr{M}^{m+1}$ 在 $m=1$ 时应为 $\scr{M}^2$，且按本节记号，上标等于其中点的个数）；此处照原书排印，存疑备案。——译者注}是\trFloerEquation{}的一个解，它以下述方式作为一个“无穷远处的\trBubble”出现。令
\begin{equation*}
\epsilon^i=\epsilon^{d+1}_{n+1}|(\{r^i\}\times Z^+)\colon Z^+\longrightarrow\scr{S}^{d+1}_{r^i}
\end{equation*}
为 $\scr{S}^{d+1}_{r^i}$ 的第 $(n{+}1)$ 个出向\trStriplikeEnd。于是存在 $\sigma^i\to+\infty$，使得\trRescaling{}映射 $\tilde{u}^i(s,t)=u^i(\epsilon^i(s+\sigma^i,t))$ 在紧子集上\trConvergence{}到 $u_2$。这个陈述之所以有意义，是因为诸 $\tilde{u}^i$ 是\trFloerEquation{}的部分解（这由加于任何\trPerturbationDatum{}之上的条件\eqref{eq:perturbation-compatibility} 推出），它们定义在一些当 $i\to\infty$ 时\trExhausting{} $Z$ 的子集上。
\item 在剩下的\emph{稳定情形} $1<m<d$ 中，$r^i$ \trConvergence{}到 $\partial\bar{\scr{R}}^{d+1}$ 中的点 $(r_1,r_2)$。按\trTopology{}的定义，这意味着存在\trConvergence{}序列 $r^i_k\to r_k$ 与一列\trGluingParameter{}$\rho^i\to 0$，使得对大的 $i$ 有 $r^i=\gamma^T(\rho^i,r^i_1,r^i_2)$。把与之相关联的\trPointedDisc{}记为
\begin{equation*}
S^i_1=\scr{S}^{d+2-m}_{r^i_1},\qquad S^i_2=\scr{S}^{m+1}_{r^i_2}
\end{equation*}
便知 $S^i$ 是把 $S^i_1$ 与 $S^i_2$ 按\eqref{eq:striplike-gluing} 那样\trGluing{}起来的结果，其\trGluingLength{}为 $l^i=-\log(-\rho^i)/\pi$。把 $u^i$ 限制到 $S^i$ 由\trGluing{}继承下来的各块上，将给出 $S_1$、$S_2$ 上的部分定义的映射序列，它们在紧子集上\trConvergence{}到 $u_1$、$u_2$。这之所以有意义，是因为\trPerturbationDatum{}选取中的\trConsistency：它要求用在 $S^i$ 与 $S^i_k$ 上的\trPerturbationDatum{}互相关联。
\end{itemize}
